\usepackage[hyperfootnotes=false]{hyperref}
\usepackage{pdfpages}
\usepackage{fontspec}
\usepackage{fancyhdr}
\usepackage[compact, bottomtitles]{titlesec}
\usepackage{alphalph}
\usepackage[para, perpage]{footmisc}
\usepackage{footnpag}
\usepackage{graphicx}
\usepackage{svg}
\usepackage{ragged2e}
\usepackage{titletoc}
\usepackage{ifthen}
\usepackage[toc]{multitoc}
\usepackage{etoolbox}
\usepackage{alphalph}
\usepackage{paracol}
\usepackage{comment}
\usepackage{titlesec}
\usepackage{xcolor}
%\usepackage[x-1a]{pdfx}
%\usepackage{bidi}
\usepackage{multicol}
\usepackage{wrapfig}
\usepackage{xparse}
\usepackage{needspace}

% A chapter title must not be left alone at the bottom of a page. multicols refuses to start unless
% \premulticols (50pt) is free, so a title with less room than that below it is pushed away from its
% text. \chapter therefore starts a new page when less than \chapterneedspace is left: roughly the title
% (about 36pt) + \multicolsep (12pt) + \premulticols (50pt). Raise it if a title still gets stranded;
% set it to 0pt to switch the check off.
\newlength{\chapterneedspace}
\setlength{\chapterneedspace}{100pt}

% Space above and below a chapter title. It is chosen inside \chapter, *after* the \needspace page-break
% decision, so that a title at the top of a page gets no space above it whether it is the first chapter
% or was pushed to a new page by \needspace. Elsewhere on the page it is separated from the text above.
% (\pagetotal is 0pt at the top of a page.)
\makeatletter
\newcommand{\biblechapter@spacing}{%
  \ifdim\pagetotal<1pt
    \titlespacing*{\chapter}{0pt}{-6mm}{0pt}%
  \else
    \titlespacing*{\chapter}{0pt}{20pt}{0pt}%
  \fi
}
\makeatother

% Do not start a chapter on a new page: extbook's \chapter begins with \cleardoublepage (\clearpage with
% openany) and \thispagestyle{plain}; the latter would also strip the header from the page the chapter
% starts on. Patching \chapter itself leaves \newpage/\clearpage working everywhere else. To get the
% default behaviour back for a book, delete this block.
\makeatletter
\patchcmd{\chapter}{\if@openright\cleardoublepage\else\clearpage\fi}{\par\needspace{\chapterneedspace}\biblechapter@spacing}{}{\PackageWarning{preamble}{could not patch \string\chapter\space (page break)}}
\patchcmd{\chapter}{\thispagestyle{plain}}{}{}{\PackageWarning{preamble}{could not patch \string\chapter\space (page style)}}
\makeatother

\extrafloats{100}
\maxdeadcycles=200
\renewcommand{\thefootnote}{\alphalph{\value{footnote}}}
\renewcommand*{\multicolumntoc}{2}
\patchcmd{\makefootnoteparagraph}
   {\columnwidth}{\textwidth}
   {\typeout{Success}}{\typeout{patch failed}}

\setmainfont{BibleRoman}
            [UprightFont = *,
             BoldFont = *Bold,
             ItalicFont = *Italic,
             Ligatures=TeX]

\makeatletter             
\newlength\fake@f
\newlength\fake@c
\def\fakesc#1{%
  \begingroup%
  \xdef\fake@name{\csname\curr@fontshape/\f@size\endcsname}%
  \fontsize{\fontdimen8\fake@name}{\baselineskip}\selectfont%
  \uppercase{#1}%
  \endgroup%
}
\makeatother
\newcommand\fauxsc[1]{\fauxschelper#1 \relax\relax}
\def\fauxschelper#1 #2\relax{%
  \fauxschelphelp#1\relax\relax%
  \if\relax#2\relax\else\ \fauxschelper#2\relax\fi%
}
\def\Hscale{.83}\def\Vscale{.72}\def\Cscale{1.00}
\def\fauxschelphelp#1#2\relax{%
  \ifnum`#1=\lccode`#1\relax\scalebox{\Hscale}[\Vscale]{\char\uccode`#1}\else%
    \scalebox{\Cscale}[1]{#1}\fi%
  \ifx\relax#2\relax\else\fauxschelphelp#2\relax\fi}

\renewcommand{\textsc}[1]{{\fauxsc{#1}}}
\renewcommand{\emph}[1]{{\textit{#1}}}

\pagestyle{fancy}
\renewcommand{\sectionmark}[1]{\markright{\thesection~- ~#1}}
\renewcommand{\chaptermark}[1]{\markboth{\chaptername~\thechapter~-~ #1}{}}
\pagestyle{empty}

\titleformat{\chapter}[display]
  {\normalfont\bfseries}{}{0pt}{\Huge}

\fancyhf{}

\fancypagestyle{bible}{%
    \fancyhead{} %Clean headers
    \fancyhead[L]{\leftmark \ \thesection}
    \fancyhead[R]{\thepage}
    \fancyhead[RO]{\leftmark \ \thesection}
    \fancyhead[LO]{\thepage}
}

\fancypagestyle{plain}{%
  \fancyhead{}
  \fancyhead[R, RO]{\thepage}
}

\renewcommand{\chaptermark}[1]{\markboth{\thechapter. {\slshape{##1}}}{}}
\renewcommand{\headrulewidth}{0pt}
% \renewcommand*{\thefootnote}{\alph{footnote}}
\titleformat{\chapter}[hang]{\normalfont\bfseries}{}{0pt}{\Huge}
\titleformat{\section}[wrap]{\bfseries\huge}{}{0ex}{}[]
\titleformat{\subsection}[hang]{\bfseries}{}{0ex}{}[]
\newcommand{\bibleverse}[1]{{\bfseries{#1}}}
\newcommand{\wordsOfJesus}[1]{{#1}}
\renewcommand{\thesection}{\arabic{section}}
\renewcommand{\chaptermark}[1]{\markboth{\MakeUppercase{#1}}{}}
\setcounter{secnumdepth}{1}
\setcounter{tocdepth}{0}
\global\let\endtitlepage\relax
\makeatletter

\titlecontents{chapter}% <section-type>
[0pt]% <left>
{\bfseries\small}% <above-code>
{\small\thecontentslabel \quad}%<numbered-entry-format>
{}% <numberless-entry-format>
{\small\mdseries\titlerule*[0.75em]{.}\bfseries\contentspage}

\defaultfontfeatures{Scale=MatchLowercase} %so that different fonts have same xheight
\newfontfamily\hebrewfont[Script=Hebrew]{Frank Ruehl CLM}
\newfontfamily\greekfont[Script=Greek]{Cardo}
\newcommand{\hebrew}[1]{{\hebrewfont{#1}}}
\newcommand{\greek}[1]{{\greekfont{#1}}}

\newcommand{\biblebeginparacol}{\begin{paracol}{2}}
\newcommand{\bibleendparacol}{\end{paracol}}

\let\oldsection\section
\let\oldchapter\chapter

\newlength{\sectionnumberwidth}
\newlength{\sectionnumberheight}
\newlength{\sectiontextwidth}

% Section numbers (the numbers of the chapters of a Bible book). The generated .tex has, for each one,
%     \hypertarget{...}{% \section{N}\label{...}}          <- \section{N} reaches \biblechapter below
%     \subsection{...}                                      <- optional headings
%     \sectionnumber \bibleverse{1} ...                     <- first paragraph of the section
% \biblechapter (called by \section for numeric titles) sets the \section counter to N. \sectionnumber
% always prints the current value of that counter (\thesection, never the \chapter number): it puts it in
% the left margin of the paragraph that starts there and indents the first two lines around it.
% \hangindent/\hangafter are set \global so the effect cannot be lost to a closing group; TeX resets
% them itself at the end of the paragraph.
\newcommand{\biblechapter}[1]{%
  \par
  % \biblechapter replaces \section for numeric titles, so it must do the counting itself (the page
  % header prints \thesection). Step to exactly N so the counter follows the number in the text even
  % if a chapter is skipped, and so \label gets the right \@currentlabel and anchor.
  \setcounter{section}{\numexpr#1-1\relax}%
  \refstepcounter{section}%
}

\newcommand{\sectionnumber}{%
  \ifnum\value{section}>\z@
    \noindent
    \settowidth{\sectionnumberwidth}{\textbf{\Huge \thesection}}%
    \addtolength{\sectionnumberwidth}{.5em}%
    \global\hangindent=\sectionnumberwidth
    \global\hangafter=-2
    % Line 1 itself starts at \hangindent, so the box is pushed back to the margin with \llap.
    % \smash keeps the large number from changing the line height; it is lowered to sit across both lines.
    \llap{\hbox to \sectionnumberwidth{\smash{\raisebox{-0.7\baselineskip}{\textbf{\Huge \thesection}}}\hss}}%
  \fi
}

\ExplSyntaxOn

\RenewDocumentCommand{\section}{m}
 {
   \regex_match:nnTF { \A \d+ \Z } { #1 }
     { \biblechapter{#1} }
     { \oldsection{#1} }
 }

\ExplSyntaxOff